\subsection{喷射流形}

12.3.1. 设 \(c = (\mathrm{U}, \varphi, \mathrm{E})\) 与 \(c' = (\mathrm{V}, \psi, \mathrm{F})\) 分别为 X 和 Y 的图表。映射 \(\varphi\) 和 \(\psi\) 通过结构迁移定义了一个双射 \(\pi\)，从 \(\mathbf{J}^k (\mathbf{U},\mathbf{V})\) 到

\[
\mathrm{J} ^ {k} (\varphi (\mathrm{U}), \psi (\mathrm{V})) = \varphi (\mathrm{U}) \times \psi (\mathrm{V}) \times \mathrm{Q} ^ {k} (\mathrm{E}, \mathrm{F}),
\]

参见 12.2.2。若置 \(\mathbf{W} = \mathbf{J}^{k}(\mathbf{U},\mathbf{V})\)，\(\mathbf{G} = \mathbf{E}\times \mathbf{F}\times \mathbf{Q}^{k}(\mathbf{E},\mathbf{F})\)，则三元组 \((\mathbf{W},\pi ,\mathbf{G})\) 是 \(\mathbf{J}^k (\mathbf{X},\mathbf{Y})\) 的一个图。由此得到的图组成一个 \(\mathbf{C}^{r - k}\)-图册，并由此赋予 \(\mathbf{J}^k (\mathbf{X},\mathbf{Y})\) 一个类为 \(\mathbf{C}^{r - k}\) 的 K-流形结构。\footnote{当 \(k=r\)（这仅在 \(K=\mathbf{R}\) 时可能）时，赋予 \(\mathbf{J}^{k}(\mathbf{X},\mathbf{Y})\) 一个拓扑流形结构（参见记号与约定），下文考虑的态射和纤维化应理解为纯粹拓扑意义下的态射和纤维化（参见 § 6 的脚注）。}

若 \((x, y) \in X \times Y\)，则 \(J_{x}^{k}(X, Y)\)、\(J^{k}(X, Y)_{y}\) 和 \(J_{x}^{k}(X, Y)_{y}\) 是 \(J^{k}(X, Y)\) 的闭子流形。

若 X 和 Y 是纯的（相应地有限维、相应地分离、相应地连通），则 \(J^{k}(X, Y)\) 也具有相应性质。

若 U（相应地 V）是 X（相应地 Y）中的开集，则 \(J^{k}(U, V)\) 是 \(J^{k}(X, Y)\) 的开子流形，参见 12.1.3。

12.3.3. 映射 \(s: \mathbf{J}^{k}(\mathbf{X}, \mathbf{Y}) \to \mathbf{X}, b: \mathbf{J}^{k}(\mathbf{X}, \mathbf{Y}) \to \mathbf{Y}\) 和 \((s, b): \mathbf{J}^{k}(\mathbf{X}, \mathbf{Y}) \to \mathbf{X} \times \mathbf{Y}\) 是类 \(\mathbf{C}^{r - k}\) 的纤维化（6.1.1）。当 \(k = 0\) 时，\((s, b)\) 是一个同构。

12.3.4. 将 \(T_{X}\) 和 \(T_{Y}\) 记为向量丛 \(\mathrm{pr}_{1}^{*}\mathrm{T}(\mathrm{X})\) 和 \(\mathrm{pr}_{2}^{*}\mathrm{T}(\mathrm{Y})\)，它们定义在 \(X \times Y\) 上，并令 \(\mathcal{L}(T_{X}; T_{Y})\) 为 \(T_{X}\) 到 \(T_{Y}\) 的同态丛（7.7.3）；若 \((x, y) \in X \times Y\)，则有

\[
\mathcal {L} (\mathrm{T} _ {\mathrm{X}}; \mathrm{T} _ {\mathrm{Y}}) _ {(x, y)} = \mathcal {L} (\mathrm{T} _ {x} (\mathrm{X}); \mathrm{T} _ {y} (\mathrm{Y})).
\]

当 \(k = 1\) 时，由此定义的映射

\[
\mathbf {T} \colon \mathbf {J} ^ {1} (\mathrm{X}, \mathrm{Y}) \rightarrow \mathcal {L} (\mathrm{T} _ {\mathrm{X}}; \mathrm{T} _ {\mathrm{Y}})
\]

是类 \(C^{r-1}\) 的流形同构。

当 X = K 时，此同构将 \(J_{0}^{1}(K, Y)\) 与 T(Y) 等同；当 Y = K 时，将 \(J^{1}(X, K)_{0}\) 与 T(X) 的对偶 \(T'(X)\) 等同。

12.3.5. 设 \(k'\) 是满足 \(0 \leqslant k' \leqslant k\) 的整数。映射

\[
r ^ {k, k ^ {\prime}}: \mathrm{J} ^ {k} (\mathrm{X}, \mathrm{Y}) \rightarrow \mathrm{J} ^ {k ^ {\prime}} (\mathrm{X}, \mathrm{Y})\tag{cf. 12.1.5}
\]

是类 \(C^{r-k}\) 的纤维化。

12.3.6. 设 Z 是类 \(\mathbf{C}^r\) 的流形，并令 T 为形如

\[
(j, j ^ {\prime}) \in \mathrm{J} ^ {k} (\mathrm{X}, \mathrm{Y}) \times \mathrm{J} ^ {k} (\mathrm{Y}, \mathrm{Z})
\]

满足 \(b(j) = s(j')\)。于是 T 是 \(J^k(X, Y) \times J^k(Y, Z)\) 的子流形，且映射 \((j, j') \mapsto j' \circ j\) 是类 \(C^{r-k}\) 的、从 T 到 \(J^k(X, Z)\) 的态射。

12.3.7. 若 \(f: \mathbf{X} \to \mathbf{Y}\) 是类 \(\mathbf{C}^r\) 的态射，则映射 \(x \mapsto j_x^k(f)\) 是一个态射 \(j^k(f)\)，类为 \(\mathbf{C}^{r-k}\)，从 \(\mathbf{X}\) 到 \(\mathbf{J}^k(\mathbf{X}, \mathbf{Y})\)。

12.3.8. 令 \(k'\) 和 \(k''\) 为和为 \(k\) 的正整数。令 \(x \in \mathbf{X}\)，令 \(\mathbf{U}\) 为 \(x\) 的一个开邻域，并令 \(f: \mathbf{U} \to \mathbf{Y}\) 是类为 \(\mathbf{C}^{r}\) 的态射。映射

\[
x \mapsto j _ {x} ^ {k ^ {\prime}} (f) \colon \mathrm{U} \to \mathrm{J} ^ {k ^ {\prime}} (\mathrm{X}, \mathrm{Y})
\]

属于类 \(\mathbf{C}^{r - k'}\)，且其 \(k''\) 阶喷射在 \(x\) 处仅取决于 \(j_x^k (f)\)。由此得到一个典范映射

\[
\alpha \colon \mathrm{J} ^ {k} (\mathrm{X}, \mathrm{Y}) \rightarrow \mathrm{J} ^ {k ^ {\prime \prime}} (\mathrm{X}, \mathrm{J} ^ {k ^ {\prime}} (\mathrm{X}, \mathrm{Y}))
\]

它属于类 \(C^{r-k}\)；若 K 的特征为零，则它是一个嵌入。

12.3.9. 设 \(\mathbf{X}'\)、\(\mathbf{Y}'\) 是类 \(\mathbf{C}^r\) 的流形，\(f: \mathbf{X} \to \mathbf{X}'\) 和 \(g: \mathbf{Y}' \to \mathbf{Y}\) 是态射，且 \((x, y') \in \mathbf{X} \times \mathbf{Y}'\)，\(x' = f(x)\)，\(y = g(y')\)。若 \(u \in \mathrm{J}_x^k(\mathbf{X}', \mathbf{Y}')_{y'}\)，置

\[
\mathrm{J} _ {x ^ {\prime}} ^ {k} (f, g) _ {y} (u) = j _ {y ^ {\prime}} ^ {k} (g) \circ u \circ j _ {x} ^ {k} (f).
\]

于是得到一个映射

\[
\mathrm{J} ^ {k} (f, g) \colon \mathrm{J} ^ {k} (\mathrm{X} ^ {\prime}, \mathrm{Y} ^ {\prime}) \times_ {\mathrm{X} ^ {\prime}} \mathrm{X} \rightarrow \mathrm{J} ^ {k} (\mathrm{X}, \mathrm{Y})
\]

它属于类 \(C^{r-k}\)。

12.3.10. 设 A 是 X 的一个紧子集。令 \(\mathcal{C}^{k}(X;Y)\) 为从 X 到 Y 的类 \(C^{k}\) 映射集合。对于每个 \(f\in\mathcal{C}^{k}(X;Y)\)，映射 \(j^{k}(f)|A\) 属于 \(\mathcal{C}(A;J^{k}(X,Y))\)，即从 A 到 \(J^{k}(X,Y)\) 的连续映射集合。由此定义了一个映射 \(\lambda\)，从 \(\mathcal{C}^{k}(X;Y)\) 到 \(\mathcal{C}(A;J^{k}(X,Y))\)。\(\lambda\) 下，定义在 \(\mathcal{C}(A;J^{k}(X,Y))\) 上的紧收敛拓扑的逆像（TG, X, § 3, 定义 1）是 \(\mathcal{C}^{k}(X;Y)\) 上的一个拓扑；称为在 A 上的 \(C^{k}\)-一致收敛拓扑。

